% !PRINTED 600
\weberpage{600}

Under the assumption, then, that Theorem IV is correct for equations of degree $\mu$, its correctness follows for equations of degree $\mu m$; and since it holds for equations of prime degree, it is proved in general.

\begin{center}
\S.\,177.

\gesperrt{Simplicity of the alternating group.}
\end{center}

We have seen earlier (\S.\,149) that if we regard the coefficients of an equation of degree $n$ as independent variables, and the field of all rational functions of these coefficients as the domain of rationality, the \gesperrt{Galois} group of the equation is the symmetric group. In the symmetric group there is always contained a normal divisor of index 2, the \gesperrt{alternating group}, to which the group of the equation reduces when the square root of the discriminant is adjoined. For four figures the alternating group, which consists of the identical permutation, eight three-term cycles, and three pairs of transpositions, has the normal divisor of index 3:

\[
1, \quad (0,\,1)(2,\,3), \quad (0,\,2)(1,\,3), \quad (0,\,3)(1,\,2).
\]

This group has three different normal divisors of index 2, of which we prefer one, 1, $(0,\,1)(2,\,3)$, of which in turn the unit group is a normal divisor of index 2. The group of the 24 permutations of four figures is therefore metacyclic, and on this rests every method of solution of the biquadratic equation (\S.\,160, 161).

We now wish to prove that if $n$ is greater than 4, the alternating group has no normal divisors at all besides the unit group, or, in the designation introduced earlier, is \gesperrt{simple}. From this it then follows that the condition which we found to be necessary for the algebraic solvability of an equation is not fulfilled for equations of higher than the fourth degree whose group is the symmetric or the alternating group, and that therefore \gesperrt{equations of higher than the fourth degree, so long as the coefficients are independent variables, are no longer algebraically solvable.}
